\chapter{Coefficient-space operators}
\label{sec:ofdg-implementation}
\label{app:reproducibility}

\section{Scope and geometric assumptions}

The implementation is a small compiled MFEM library.  It assumes that the
finite element space is nonempty and discontinuous, that $k\geq 1$, and that
all elements have the same order and dimension.  Geometry and local degree-of-
freedom counts may differ.  Element maps may be affine or static polynomial curvilinear, and the mesh
must remain unchanged during the lifetime of the OFDG object.  The space and cache-lifetime conditions are caller responsibilities. Native NURBS geometry is explicitly rejected before the unsupported face path is entered.
The public header exposes the OFDG constructors, stabilization, and decay only.
KXRCF is an independent optional method and is included explicitly when the
activation mask is required.

Different elements need not have the same Jacobian.  They may have different
sizes, reference geometries, and local coefficient counts, and may be
translated, rotated, or affinely skewed.  An immutable repository constructs
the required projection and derivative operators once per finite-element
signature.  Non-affine elements additionally own physical projection and
projected-derivative matrices assembled with geometry-dependent quadrature.
Native NURBS geometry is rejected: the pinned MFEM implementation does not
support the two-sided NURBS face transformations required here.  Projecting such
a mesh to polynomial nodes changes the rational geometry approximately and is
therefore a distinct supported case.

More explicitly, MFEM describes a physical element by a map from a standard
reference element.  The affine assumption means that this map has the form
\begin{equation}
  \boldsymbol{x}=T_K(\boldsymbol{\xi})
  =\boldsymbol{b}_K+J_K\boldsymbol{\xi},
  \label{eq:ofdg-affine-map}
\end{equation}
where $\boldsymbol{\xi}$ is a reference coordinate, $\boldsymbol{b}_K$ is a
translation, and the matrix $J_K$ is constant within $K$.  The matrix may differ
from one element to another.  Its scaling, rotation, and shear describe the size,
orientation, and skew of the physical element.  The affine path therefore uses a constant Jacobian on each individual element,
not one globally identical Jacobian.  For curved maps the volume is integrated
and the damping length is $h_K=(|K|/|\widehat K|)^{1/d}$, agreeing with the
existing affine determinant scale.

\section{Projection matrices}
The physical change of variables and projection derivation are given in
\cref{sec:curvilinear-mathematics}. Here we specify the stored remainder matrices.

Let $\{\phi_i\}_{i=1}^{N_k}$ be the local degree-$k$ basis and
$\{\psi_a^{(r)}\}_{a=1}^{N_r}$ the degree-$r$ basis.  The code assembles
\begin{align}
  (M_k)_{ij}&=(\phi_i,\phi_j)_K,
  & (M_r)_{ab}&=(\psi_a^{(r)},\psi_b^{(r)})_K,\\
  (B_r)_{ai}&=(\psi_a^{(r)},\phi_i)_K.
\end{align}
The matrix representation of $I-\Pi_K^r$ in the degree-$k$ coefficient basis is
\begin{equation}
  Q_r=M_k-B_r^\mathsf{T}M_r^{-1}B_r,
  \qquad
  S_r=M_k^{-1}Q_r
     =I-M_k^{-1}B_r^\mathsf{T}M_r^{-1}B_r.
  \label{eq:ofdg-projection-matrices}
\end{equation}
Thus applying $S_r$ to a degree-$k$ coefficient vector directly returns the
coefficients of the remainder $w-\Pi_K^r w$.  The matrices $M_k$ and $M_r$ are
the degree-$k$ and degree-$r$ mass matrices, respectively; $B_r$ couples the two
bases; and $I$ is the $N_k\times N_k$ identity matrix.  The code stores $S_r$
for $0\leq r\leq k$.

Crucially, this construction never converts the solution to Legendre or other
orthogonal modal coefficients.  It represents the same projection operator in
whatever supported scalar polynomial DG basis the host framework supplies.
Consequently, a nodal basis, a modal basis, or another polynomial coefficient
basis produces the same function-space operation up to quadrature and roundoff.
This basis-agnostic realization is what makes the implementation portable: a
second framework need not add a special OFDG solution basis, provided it can
form the corresponding polynomial-space projections.

On an affine element, all three mass
and coupling matrices contain the same constant physical Jacobian factor, which
cancels in $S_r$.  Therefore affine projection matrices are reused per finite-element signature.
For curved elements every mass and coupling matrix includes the physical
Jacobian weight at each quadrature point.  The lower-order spaces are mapped
reference polynomial spaces; their physical $L^2$ projectors remain nested.

\section{Coefficient-to-coefficient derivatives}

The repeated derivatives required by~\eqref{eq:ofdg-face-jump} are represented as
polynomial coefficients rather than only as values at quadrature points.  At the
finite element interpolation points $\{\boldsymbol{\xi}_q\}_{q=1}^{N_k}$, define
\begin{equation}
  V_{qi}=\phi_i(\boldsymbol{\xi}_q),
  \qquad
  G_{qi}^{(r)}
  =\frac{\partial\phi_i}{\partial\xi_r}(\boldsymbol{\xi}_q).
\end{equation}
If $\boldsymbol{u}$ is a coefficient vector, $G^{(r)}\boldsymbol{u}$ contains
reference-derivative values at the interpolation points.  The coefficient vector
of that derivative is therefore
\begin{equation}
  \boldsymbol{u}_{\xi_r}=D_{\xi_r}\boldsymbol{u},
  \qquad
  D_{\xi_r}=V^{-1}G^{(r)}.
  \label{eq:ofdg-reference-derivative-matrix}
\end{equation}
This Vandermonde construction is kept explicit in the code.  For the affine map
$\boldsymbol{x}=T_K(\boldsymbol{\xi})$, the physical derivative matrices are
\begin{equation}
  D_{x_j}^{K}
  =\sum_{r=1}^d(J_K^{-1})_{rj}D_{\xi_r}.
  \label{eq:ofdg-physical-derivative-matrix}
\end{equation}
The inverse Jacobian $J_K^{-1}$ is read once at the reference-element centre and
one set of matrices $\{D_{x_j}^{K}\}_{j=1}^d$ is stored for every element.

For non-affine maps, the code instead assembles
\begin{equation}
  M_K D_{K,j}=G_{K,j},\qquad
  (G_{K,j})_{ab}=\int_K\phi_a\,\partial_{x_j}\phi_b\,dx.
\end{equation}
Repeated application represents successively projected physical derivatives,
not exact higher derivatives of the mapped function.  The derivative graph
uses $j=\max\{s:\alpha_s>0\}$ and parent $\alpha-e_j$ for every
nonzero multi-index. It therefore applies directions in the order stated in
\cref{eq:projected-multi-index-order}; projected mixed derivatives need not
commute. Exact physical mixed partials still commute for sufficiently smooth
functions, as distinguished in \cref{eq:curved-commutator}.
Positive, overintegrated quadrature is used for curved elements.  Triangular
face rules include all barycentric permutations to make the sampling invariant
under MPI face orientation.

Only first reference derivatives of the supplied finite-element basis are
requested from MFEM.  Repeated application of the resulting coefficient
operators generates every higher derivative needed by OFDG.  In framework
terms, the construction requires scalar polynomial DG basis values and first
derivatives, quadrature, nested lower-order polynomial spaces or equivalent
projection operators, element and face transformations, and local
degree-of-freedom access.  It therefore avoids a modal-basis requirement, but it
does not claim support for arbitrary nonpolynomial, $H(\mathrm{div})$,
$H(\mathrm{curl})$, or mixed finite-element spaces.

All multi-indices with total degree at most $k$ are organized in a directed
acyclic graph.  Each nonzero multi-index selects one parent of degree one less and
one physical derivative direction.  Walking this graph evaluates every required
derivative once.  The number of stored derivative states per element and component
is
\begin{equation}
  \sum_{l=0}^k\binom{l+d-1}{d-1}=\binom{k+d}{d}.
\end{equation}

\section{Face quadrature and cached evaluation}

MFEM remains the sole source of mesh connectivity.  The setup loops over MFEM
face numbers and obtains the two adjacent elements from MFEM's face-transformation
data.  Boundary faces are skipped.  For each
interior affine face, a quadrature rule of order $2k+1$ is selected.  Curved
faces use an overintegrated rule that also depends on the geometry order.  The following data are cached directly at the MFEM face index:
\begin{equation}
  (E_F^-)_{qi}=\phi_i(\boldsymbol{\xi}_{F,q}^-),
  \qquad
  (E_F^+)_{qi}=\phi_i(\boldsymbol{\xi}_{F,q}^+),
  \qquad
  \widehat{\omega}_q=\frac{\omega_q J_F(\xi_q)}{\sum_p\omega_p J_F(\xi_p)}.
  \label{eq:ofdg-face-evaluation-matrices}
\end{equation}
Here $J_F$ is the physical face Jacobian.  For affine faces it is constant and cancels between the
integral and $|F|$ in~\eqref{eq:ofdg-face-jump}; hence the normalized reference
weights $\widehat{\omega}_q$ are sufficient.  All derivative traces and all
components on one side of a face are evaluated together in one dense matrix
multiplication.  There is no face-pattern matching or deduplication.

If $\boldsymbol{u}_{\boldsymbol{\alpha}}^-$ and
$\boldsymbol{u}_{\boldsymbol{\alpha}}^+$ are the coefficient vectors of the
same physical derivative on the two neighboring elements, the actual quadrature
calculation is
\begin{equation}
  J_{F,c}^{(l)}
  \approx\sum_{|\boldsymbol{\alpha}|=l}\sum_{q=1}^{N_{q,F}}
  \widehat{\omega}_q
  \left|
    \bigl(E_F^-\boldsymbol{u}_{\boldsymbol{\alpha}}^-\bigr)_q
    -\bigl(E_F^+\boldsymbol{u}_{\boldsymbol{\alpha}}^+\bigr)_q
  \right|.
  \label{eq:ofdg-discrete-face-jump}
\end{equation}
The implementation therefore performs one subtraction and one absolute value for
each derivative, component, and quadrature point.  No squared-jump accumulator or
final square root is required.

\section{Runtime sequence and storage policy}

The complete filter application is specified in
\cref{sec:ofdg-update-recipe}. The public semidiscrete alternative uses the
same rates and returns the coefficient vector of
\cref{eq:ofdg-damping-operator} instead of applying exponential decay.
Affine volume means use a rule of order $2k+1$; curved elements use the geometry-dependent overintegration described in \cref{sec:curved-construction}.  The extrema used for
$A_c$ are the minimum and maximum over the same quadrature values.  In floating
point arithmetic, the implementation sets the rates of component $c$ to zero when
$A_c\leq 10^{-14}$.

There is no bounded or dynamic derivative cache.  Every element owns permanent
storage for its derivative coefficients, and all of it is rebuilt once for the
current state.  Likewise, each interior face owns its two evaluation matrices and
normalized quadrature weights.  Ignoring small container overheads, the principal
storage costs are
\begin{align}
  &\mathcal{O}\left(
    |\mathcal{T}_h|\,dN_k^2
    \right)
  &&\text{for physical derivative operators},\\
  &\mathcal{O}\left(
    |\mathcal{T}_h|\,mN_k\binom{k+d}{d}
    \right)
  &&\text{for derivative coefficient states},\\
  &\mathcal{O}\left(
    \sum_{F\in\mathcal{F}_h^{\mathrm{i}}}N_{q,F}(2N_k+1)
    \right)
  &&\text{for face evaluation data}.
\end{align}
This policy uses more predictable memory than an evicting cache and removes cache
lookups and replacement decisions from the face loop.

An optional active-element mask can restrict stabilization or decay to selected
elements.  A local face is omitted before constructing its geometry when neither
neighbor is active, so an active
element still receives contributions from every adjacent interior face. Runtime
scratch arrays are owned by the OFDG object. Because this scratch state is
mutable and reused, concurrent calls on the same object are unsupported.


\section{Verification criteria}
\label{sec:verification-criteria}
The algebra suggests independent checks: constants and component-wise cell
means should be preserved; inactive cells should remain unchanged; and a
change of coefficient basis should represent the same filtered function.
Curved first derivatives can be compared with an independently assembled
physical projection. Recursive higher derivatives require manufactured-solution
refinement studies, since their definition alone supplies no convergence
rate. Quadrature sensitivity, geometric parameterization, and serial/MPI
agreement are separate checks of the geometric realization.

These checks distinguish consistency of an implementation from the accuracy
and robustness of a complete simulation. Their selected numerical evidence,
along with comparisons between methods, belongs in the Results chapter once
the experimental programme has been agreed.
